\itemhead{第3章 小结与复习 · 第3题（教材p.89）}{XJ-B1-C03-SU-P089-Y1112}
\label{XJ-B1-C03-SU-P089-Y1112}
{\footnotesize 湘教版教材p.89 / PDF第96页 / 第3章 小结与复习 / 第3题}
\field{原题}{根据函数的解析式或图象来探索它的性质，是数学的重要课题。函数的变化无非是变大变小，函数的单调性描述了函数的变化过程和趋势，是函数最重要的特征。如何判断一个函数的单调性？更进一步，你能用数学语言来表述吗？函数最值的几何意义如何理解？}
\field{参考答案}{单调性用区间内任意 $x_1<x_2$ 的函数值次序定义；最值 $M$ 要同时满足所有函数值不超过 $M$，并有某点实际取到 $M$。几何上为图象实际达到的最高或最低纵坐标。}
{\small \field{关联知识}{F05 单调性定义、证明与参数（教材帮p.97、98、100 / PDF106、107、109）；F06 最大值、最小值与值域（教材帮p.99、83、84 / PDF108、92、93）}
\field{方法与易错点}{用定义比较任意两值，或依据已知图象性质；含参时使指定区间被相应单调区间覆盖。 易错：用有限采样代替证明；把分离半轴合并；漏跨零点情形。}
\field{答案依据}{答案册未提供本栏目的独立对应答案；依据原题、教材概念补写参考回应。开放性总结不设唯一数值答案。}}
